% !TEX root = ../main.tex
\corrigesChapitre{1}{Transformations acide-base}
\corrigesSouples

\begin{blocLivre}{cor-1-1}
\begin{corrige}{1}{Acide ou base ?}
\demarche{Associer chaque espèce à celle qui en diffère d'un seul ion \ce{H+}.}
\begin{center}\small\renewcommand{\arraystretch}{1.25}
\begin{tabular}{@{}llll@{}}\toprule
\otTete{Espèce} & \otTete{Nature} & \otTete{Espèce conjuguée} & \otTete{Couple}\\\midrule
\ce{CH3COOH} & acide & \ce{CH3COO-} & \ce{CH3COOH}/\ce{CH3COO-}\\
\ce{NH3} & base & \ce{NH4+} & \ce{NH4+}/\ce{NH3}\\
\ce{HCOO-} & base & \ce{HCOOH} & \ce{HCOOH}/\ce{HCOO-}\\
\ce{CH3NH3+} & acide & \ce{CH3NH2} & \ce{CH3NH3+}/\ce{CH3NH2}\\
\ce{HO-} & base & \ce{H2O} & \ce{H2O}/\ce{HO-}\\
\ce{H3O+} & acide & \ce{H2O} & \ce{H3O+}/\ce{H2O}\\
\ce{HCO3-} & acide et base & \ce{CO3^2-} et \ce{CO2,H2O} & \ce{HCO3-}/\ce{CO3^2-} et \ce{CO2,H2O}/\ce{HCO3-}\\\bottomrule
\end{tabular}\end{center}
L'ion \res{\ce{HCO3-}} appartient à deux couples : c'est une espèce amphotère. L'eau, conjuguée
de \ce{H3O+} et de \ce{HO-}, l'est aussi, mais elle ne figure pas dans la liste.
\end{corrige}
\end{blocLivre}
\begin{blocLivre}{cor-1-2}
\begin{corrige}{2}{Écrire des demi-équations}
\demarche{Écrire acide = base + \ce{H+}, puis compter les atomes et les charges de chaque côté.}
\begin{questions}
\item \begin{tabular}[t]{@{}l@{\hspace{12mm}}l@{}}
\ce{HCOOH = HCOO- + H+} & \ce{NH4+ = NH3 + H+}\\[0.8mm]
\ce{H2O = HO- + H+} & \ce{CO2,H2O = HCO3- + H+}
\end{tabular}
\item Chaque fois, l'acide contient un atome d'hydrogène de plus que la base, porté par l'ion \ce{H+}.
Charges : $0=-1+1$ ; $+1=0+1$ ; $0=-1+1$ ; $0=-1+1$. Les éléments et la charge sont conservés.
\item \ce{HSO4-} et \ce{SO4^2-} ne diffèrent que d'un ion \ce{H+} : ils forment le couple
\res{\ce{HSO4-}/\ce{SO4^2-}}, de demi-équation \ce{HSO4- = SO4^2- + H+}.
\ce{H2SO4} et \ce{SO4^2-} diffèrent de deux ions \ce{H+} : ils ne sont pas conjugués.
Le couple de l'acide sulfurique est \ce{H2SO4}/\ce{HSO4-}.
\end{questions}
\end{corrige}
\end{blocLivre}
\begin{blocLivre}{cor-1-3}
\begin{corrige}{3}{Du pH à la concentration, et inversement}
\demarche{Utiliser $\pH=-\log(\hho/\cst)$ ou sa réciproque, en respectant les chiffres significatifs.}
\begin{questions}
\item a) $\pH=-\log(\num{1.0e-3})=\res{3{,}00}$\par
b) $\pH=-\log(\num{2.5e-4})=\res{3{,}60}$\par
c) $\pH=-\log(\num{6.3e-9})=\res{8{,}20}$\par
Deux chiffres significatifs pour la concentration donnent deux décimales pour le pH.
\item a) $\hho=10^{-2{,}0}\times\cst=\res{\qty{1e-2}{mol.L^{-1}}}$\par
b) $\hho=10^{-7{,}40}\times\cst=\res{\qty{4,0e-8}{mol.L^{-1}}}$\par
c) $\hho=10^{-11{,}3}\times\cst=\res{\qty{5e-12}{mol.L^{-1}}}$
\item À 25 °C, les solutions 1a, 1b et 2a sont acides ($\pH<7$) ; les solutions 1c, 2b et 2c sont
basiques ($\pH>7$). Aucune n'est neutre.
\end{questions}
\end{corrige}
\end{blocLivre}
\begin{blocLivre}{cor-1-4}
\begin{corrige}{4}{L'acide des fourmis}
\demarche{Répartir les électrons de valence en doublets, puis comparer les polarités des liaisons.}
\begin{questions}
\item Le carbone forme quatre liaisons, chaque oxygène deux liaisons et porte deux doublets non liants,
chaque hydrogène une liaison. Le groupe caractéristique est le \res{\text{groupe carboxyle } \ce{-COOH}}.
\item L'hydrogène cédé est celui de la liaison O--H : $\Delta\chi=3{,}44-2{,}20=1{,}24$, la liaison est
très polarisée. La liaison C--H est peu polarisée ($\Delta\chi=0{,}35$) : son hydrogène n'est pas cédé.
\item Dans l'ion méthanoate, l'oxygène qui a perdu l'ion \ce{H+} n'est lié qu'une fois et porte trois
doublets non liants : sa charge vaut $6-6-1=\res{-1}$.
\end{questions}
\begin{center}
\chemfig{H-@{c4a}C(=[1]@{c4b}\charge{0=\|,90=\|}{O})-[7]@{c4c}\charge{180=\|,270=\|}{O}-@{c4d}H}%
\chemmove{\node[pcGroupe,fit=(c4a)(c4b)(c4c)(c4d),inner sep=3.4pt]{};}
\hspace{18mm}
\chemfig{H-C(=[1]\charge{0=\|,90=\|}{O})-[7]\charge{0=\|,180=\|,270=\|,60:3.6pt=\cmoins}{O}}\par
\legende{Acide méthanoïque (groupe carboxyle entouré) et ion méthanoate.}
\end{center}
\end{corrige}
\end{blocLivre}
\begin{blocLivre}{cor-1-5}
\begin{corrige}{5}{Espèces amphotères}
\demarche{Une espèce amphotère apparaît comme acide dans un couple et comme base dans un autre.}
\begin{questions}
\item Couples \ce{H3O+}/\ce{H2O} et \ce{H2O}/\ce{HO-} : l'eau est la base du premier et l'acide du
second, elle est donc \res{\text{amphotère}}.
\item \ce{H3O+ = H2O + H+} et \ce{HO- + H+ = H2O}, d'où \res{\ce{H3O+ + HO- -> 2 H2O}}.
\item \res{\ce{H2PO4-}} est la base du couple \ce{H3PO4}/\ce{H2PO4-} et l'acide du couple
\ce{H2PO4-}/\ce{HPO4^2-} ; \res{\ce{HPO4^2-}} est la base de ce deuxième couple et l'acide du couple
\ce{HPO4^2-}/\ce{PO4^3-}. Ces deux ions sont amphotères ; \ce{H3PO4} n'est qu'acide et \ce{PO4^3-}
n'est que base.
\item L'ion carbonate n'appartient qu'au couple \ce{HCO3-}/\ce{CO3^2-}, comme base : il peut capter
un ion \ce{H+}, mais il n'en possède aucun à céder. Il n'est \res{\text{pas amphotère}}.
\end{questions}
\end{corrige}
\end{blocLivre}
\begin{blocLivre}{cor-1-6}
\begin{corrige}{6}{Établir l'équation d'une réaction}
\demarche{Écrire la demi-équation de l'acide dans le sens où il cède \ce{H+}, celle de la base dans
le sens où elle le capte, puis additionner.}
\begin{questions}
\item Acide \ce{CH3COOH} (couple \ce{CH3COOH}/\ce{CH3COO-}), base \ce{NH3} (couple \ce{NH4+}/\ce{NH3}) :\\
\methode{additionner deux demi-équations}\enspace$\res{\ce{CH3COOH + NH3 -> CH3COO- + NH4+}}$.
\item Acide \ce{H3O+} (couple \ce{H3O+}/\ce{H2O}), base \ce{CH3NH2} (couple \ce{CH3NH3+}/\ce{CH3NH2}) :\\
$\res{\ce{CH3NH2 + H3O+ -> CH3NH3+ + H2O}}$.
\item Acide \ce{HCO3-} (couple \ce{HCO3-}/\ce{CO3^2-}), base \ce{HO-} (couple \ce{H2O}/\ce{HO-}) :\\
$\res{\ce{HCO3- + HO- -> CO3^2- + H2O}}$.
\item Acide \ce{HCOOH} (couple \ce{HCOOH}/\ce{HCOO-}), base \ce{CO3^2-} (couple \ce{HCO3-}/\ce{CO3^2-}) :\\
$\res{\ce{HCOOH + CO3^2- -> HCOO- + HCO3-}}$.
\item Non. Dans la réaction 2, l'eau est la base conjuguée de \ce{H3O+} : elle apparaît quand \ce{H3O+}
cède son ion \ce{H+}. Dans la réaction 3, l'eau est l'acide conjugué de \ce{HO-} : elle apparaît quand
\ce{HO-} capte un ion \ce{H+}. L'eau joue les deux rôles : elle est amphotère.
\end{questions}
\end{corrige}
\end{blocLivre}
\begin{blocLivre}{cor-1-7}
\begin{corrige}{7}{Souffler dans une solution de bleu de bromothymol}
\demarche{Traduire chaque teinte en domaine de pH, puis relier le pH à la formation ou à la consommation d'ions \ce{H3O+}.}
\begin{questions}
\item La solution passe du vert au jaune : son pH devient inférieur à 6,0, il \res{\text{diminue}}.
La concentration en ion oxonium \res{\text{augmente}}.
\item Couples \ce{CO2,H2O}/\ce{HCO3-} et \ce{H3O+}/\ce{H2O} :
\[\res{\ce{CO2,H2O(aq) + H2O(l) <=> HCO3^-(aq) + H3O^+(aq)}}.\]
Cette réaction forme des ions oxonium : le pH diminue et le BBT vire au jaune.
\item Les ions hydroxyde réagissent avec les ions oxonium (\ce{H3O+ + HO- -> 2 H2O}) et avec le
dioxyde de carbone dissous :
\[\res{\ce{CO2,H2O(aq) + HO^-(aq) -> HCO3^-(aq) + H2O(l)}}.\]
Une fois \ce{H3O+} et \ce{CO2,H2O} consommés, les ions \ce{HO-} ajoutés en excès rendent la solution
basique : son pH dépasse 7,6 et le BBT devient bleu.
\item L'eau, base du couple \ce{H3O+}/\ce{H2O} et acide du couple \ce{H2O}/\ce{HO-}, et l'ion
hydrogénocarbonate, base du couple \ce{CO2,H2O}/\ce{HCO3-} et acide du couple \ce{HCO3-}/\ce{CO3^2-}.
\end{questions}
\end{corrige}
\end{blocLivre}
\begin{blocLivre}{cor-1-8}
\begin{corrige}{8}{Une amine et son ion ammonium}
\demarche{Repérer le doublet non liant de l'azote : c'est lui qui capte l'ion \ce{H+}.}
\begin{questions}
\item L'azote est lié à un carbone et à deux hydrogènes ; il porte un doublet non liant, site
\res{\text{basique}} de la molécule.
\item Ion éthylammonium : \res{\ce{CH3-CH2-NH3+}}. L'azote forme quatre liaisons et ne porte plus de
doublet non liant : sa charge vaut $5-0-4=+1$.
\item Couples \ce{CH3CH2NH3+}/\ce{CH3CH2NH2} et \ce{H2O}/\ce{HO-} :
\[\res{\ce{CH3CH2NH2(aq) + H2O(l) <=> CH3CH2NH3^+(aq) + HO^-(aq)}}.\]
\item $\hho=10^{-11{,}8}\times\cst=\res{\qty{2e-12}{mol.L^{-1}}}$ (une décimale pour le pH, un chiffre
significatif). Le pH est supérieur à 7 : la solution est basique, ce qui est cohérent avec la formation
d'ions hydroxyde à la question 3.
\end{questions}
\begin{center}
\chemfig{H-C(-[2]H)(-[6]H)-C(-[2]H)(-[6]H)-\charge{90=\|}{N}(-[6]H)-H}
\hspace{14mm}
\chemfig{H-C(-[2]H)(-[6]H)-C(-[2]H)(-[6]H)-\charge{45:4.2pt=\cplus}{N}(-[2]H)(-[6]H)-H}\par
\legende{Éthylamine et ion éthylammonium.}
\end{center}
\end{corrige}
\end{blocLivre}
\begin{blocLivre}{cor-1-9}
\begin{corrige}{9}{L'acide lactique}
\demarche{Seul l'hydrogène du groupe carboxyle est cédé ; le groupe hydroxyle est conservé.}
\begin{questions}
\item La molécule porte deux groupes caractéristiques : le groupe \res{\text{hydroxyle } \ce{-OH}},
de la famille des alcools, et le groupe \res{\text{carboxyle } \ce{-COOH}}, de la famille des acides carboxyliques.
\item Couple \ce{CH3-CH(OH)-COOH}/\ce{CH3-CH(OH)-COO-} ; demi-équation
\ce{CH3CH(OH)COOH = CH3CH(OH)COO- + H+}. L'ion lactate a pour formule \res{\ce{CH3-CH(OH)-COO-}}.
\item Voir le schéma ci-dessous : l'oxygène porteur de la charge est lié une seule fois et entouré de
trois doublets non liants.
\item Couples \ce{CH3CH(OH)COOH}/\ce{CH3CH(OH)COO-} et \ce{CO2,H2O}/\ce{HCO3-} :
\[\res{\ce{CH3CH(OH)COOH + HCO3- -> CH3CH(OH)COO- + CO2,H2O}}.\]
Le dioxyde de carbone dissous est transporté par le sang jusqu'aux poumons, où il est éliminé
sous forme gazeuse dans l'air expiré : \ce{CO2,H2O -> CO2(g) + H2O(l)}.
\end{questions}
\begin{center}
\chemfig{H-C(-[2]H)(-[6]H)-C(-[2]H)(-[6]\charge{0=\|,180=\|}{O}-[6]H)-C(=[1]\charge{0=\|,90=\|}{O})-[7]\charge{0=\|,180=\|,270=\|,60:3.6pt=\cmoins}{O}}\par
\legende{Ion lactate.}
\end{center}
\end{corrige}
\end{blocLivre}
\begin{blocLivre}{cor-1-10}
\begin{corrige}{10}{Une eau minérale riche en ions hydrogénocarbonate}
\demarche{Convertir la concentration en masse en concentration en quantité de matière, puis exploiter la stœchiométrie 1 : 1.}
\begin{questions}
\item $M(\ce{HCO3-})=1{,}0+12{,}0+3\times16{,}0=\qty{61,0}{g.mol^{-1}}$, d'où
\[c=\frac{c_\mathrm{m}}{M}=\frac{\qty{3,050}{g.L^{-1}}}{\qty{61,0}{g.mol^{-1}}}=\res{\qty{5,00e-2}{mol.L^{-1}}}.\]
\item \ce{HCO3-} est la base du couple \ce{CO2,H2O}/\ce{HCO3-} et l'acide du couple \ce{HCO3-}/\ce{CO3^2-} :
il est amphotère.
\item $\res{\ce{HCO3^-(aq) + H3O^+(aq) -> CO2,H2O(aq) + H2O(l)}}$. Le dioxyde de carbone dissous se dégage
en partie sous forme gazeuse : \ce{CO2,H2O -> CO2(g) + H2O(l)}, d'où les éructations.
\item $n(\ce{HCO3-})=c\times V=\num{5.00e-2}\times\num{0.200}=\qty{1,00e-2}{mol}$. Une mole d'ions
\ce{HCO3-} consomme une mole d'ions \ce{H3O+} : au plus \res{\qty{1,00e-2}{mol}} d'ions oxonium. Il se
forme au plus \qty{1,00e-2}{mol} de dioxyde de carbone, soit
$V=n\times V_\mathrm{m}=\num{1.00e-2}\times\num{24.0}=\res{\qty{0,240}{L}}$.
\end{questions}
\end{corrige}
\end{blocLivre}
\begin{blocLivre}{cor-1-11}
\begin{corrige}{11}{Lire une variation de pH}
\demarche{Une variation $\Delta\pH$ correspond à un facteur $10^{\Delta\pH}$ sur la concentration en ion oxonium.}
\begin{questions}
\item Le pH diminue de $5{,}6-4{,}3=1{,}3$ : \hho\ est multipliée par $10^{1{,}3}=10\times10^{0{,}3}\approx\res{20}$.
\item Le rapport vaut $10^{7{,}45-7{,}35}=10^{0{,}1}\approx\res{1{,}26}$. Entre les deux limites, \hho\ ne varie
que de 26 \% : l'organisme maintient cette concentration dans un intervalle étroit. Une petite variation de pH
traduit pourtant une variation notable de concentration.
\item $\pH'=-\log\left(\dfrac{\hho}{5\cst}\right)=\pH+\log5=2{,}7+0{,}70=\res{3{,}4}$.
\item Augmenter le pH de 0,5, c'est diviser \hho\ par $10^{0{,}5}\approx3{,}2$ : il faut diluer environ
\res{3{,}2 \text{ fois}}.
\end{questions}
\end{corrige}
\end{blocLivre}
\begin{blocLivre}{cor-1-12}
\begin{corrige}{12}{L'odeur du poisson}
\demarche{Une amine capte un ion \ce{H+} sur le doublet non liant de son azote et devient un ion, non volatil.}
\begin{questions}
\item L'azote est lié à trois atomes de carbone et porte un \res{\text{doublet non liant}} : c'est le site
basique de la molécule (schéma ci-dessous).
\item Couple \res{\ce{(CH3)3NH+}/\ce{(CH3)3N}}. Dans l'ion triméthylammonium, l'azote forme quatre liaisons
et ne porte plus de doublet non liant : sa charge vaut $5-0-4=+1$.
\item Couples \ce{(CH3)3NH+}/\ce{(CH3)3N} et \ce{CH3COOH}/\ce{CH3COO-} :
\[\res{\ce{(CH3)3N + CH3COOH -> (CH3)3NH+ + CH3COO-}}.\]
\item La triméthylamine, volatile et odorante, est transformée en ions triméthylammonium qui restent
dans la phase aqueuse : moins d'amine passe dans l'air, l'odeur est atténuée.
\item $\res{\ce{(CH3)3N + AH -> (CH3)3NH+ + A-}}$. Les ions \ce{Na+} et \ce{Cl-} ne contiennent pas
d'hydrogène : ils ne peuvent céder aucun ion \ce{H+} à l'amine, qui reste volatile.
\end{questions}
\begin{center}
\chemfig{H-C(-[2]H)(-[5]H)-\charge{90=\|}{N}(-[6]C(-[5]H)(-[7]H)-[6]H)-C(-[2]H)(-[7]H)-H}
\hspace{14mm}
\chemfig{H-C(-[2]H)(-[5]H)-\charge{45:4.2pt=\cplus}{N}(-[2]H)(-[6]C(-[5]H)(-[7]H)-[6]H)-C(-[2]H)(-[7]H)-H}\par
\legende{Triméthylamine et ion triméthylammonium.}
\end{center}
\end{corrige}
\end{blocLivre}
\begin{blocLivre}{cor-1-13}
\begin{corrige}{13}{Un ballon qui se gonfle tout seul}
\demarche{Calculer les quantités initiales, identifier le réactif limitant, puis convertir la quantité de gaz en volume.}
\begin{questions}
\item Couples \ce{CH3COOH}/\ce{CH3COO-} et \ce{CO2,H2O}/\ce{HCO3-} :
\[\res{\ce{CH3COOH(aq) + HCO3^-(aq) -> CH3COO^-(aq) + CO2,H2O(aq)}}.\]
Le dioxyde de carbone formé se dégage : \ce{CO2,H2O(aq) -> CO2(g) + H2O(l)}. L'effervescence est due au \res{\ce{CO2}}.
\item Un litre de vinaigre a une masse de $\qty{1,01}{g.mL^{-1}}\times\qty{1000}{mL}=\qty{1010}{g}$ et contient
$0{,}060\times\qty{1010}{g}=\qty{60,6}{g}$ d'acide, soit $\frac{60{,}6}{60{,}0}=\qty{1,01}{mol}$ :
$c=\res{\qty{1,01}{mol.L^{-1}}}\approx\qty{1,0}{mol.L^{-1}}$.
\item $n(\ce{CH3COOH})=\num{1.01}\times\num{0.0500}=\qty{5,05e-2}{mol}$ et
$n(\ce{HCO3-})=\frac{5{,}00}{84{,}0}=\qty{5,95e-2}{mol}$. Les nombres stœchiométriques sont égaux :
l'acide éthanoïque est le \res{\text{réactif limitant}} et $x_\mathrm{max}=\res{\qty{5,05e-2}{mol}}$.
\item Il se forme au plus $n(\ce{CO2})=x_\mathrm{max}$, soit $V=\num{5.05e-2}\times\num{24.0}=\qty{1,21}{L}$.
Une boule de 15 cm de diamètre a pour volume
\[\frac43\pi\times(\qty{7,5}{cm})^3=\qty{1,77e3}{cm^3}=\qty{1,77}{L}>\qty{1,21}{L}:\]
\res{\text{le ballon n'atteint pas 15 cm}}. Son diamètre maximal vaut
$D=\left(\frac{6V}{\pi}\right)^{1/3}=\left(\frac{6\times\num{1212}}{\pi}\right)^{1/3}\approx\qty{13,2}{cm}$.
\item Il faut $n=\frac{1{,}77}{24{,}0}=\qty{7,36e-2}{mol}$ de dioxyde de carbone, donc autant d'acide et
d'hydrogénocarbonate :
\[V_\text{vinaigre}\geq\frac{\num{7.36e-2}}{1{,}01}\approx\qty{73}{mL}\qquad\text{et}\qquad
m\geq\num{7.36e-2}\times84{,}0=\qty{6,19}{g}.\]
Par exemple, \res{80 \text{ mL de vinaigre et } 7 \text{ g d'hydrogénocarbonate}}, pour garder une petite marge.
\end{questions}
\end{corrige}
\end{blocLivre}
\begin{blocLivre}{cor-1-14}
\begin{corrige}{14}{L'acidification des océans}
\demarche{Relier la dissolution du dioxyde de carbone à la formation d'ions \ce{H3O+} et à la consommation d'ions carbonate.}
\begin{questions}
\item \ce{CO2,H2O}/\ce{HCO3-} et \ce{HCO3-}/\ce{CO3^2-} : l'ion hydrogénocarbonate est \res{\text{amphotère}}.
\item Couples \ce{CO2,H2O}/\ce{HCO3-} et \ce{H3O+}/\ce{H2O} :
\[\res{\ce{CO2,H2O(aq) + H2O(l) <=> HCO3^-(aq) + H3O^+(aq)}}.\]
La réaction forme des ions oxonium : \hho\ augmente, donc le pH diminue.
\item \ce{CO2,H2O = HCO3- + H+} et \ce{CO3^2- + H+ = HCO3-}, d'où
\[\res{\ce{CO2,H2O + CO3^2- <=> 2 HCO3-}}.\]
Plus il y a de dioxyde de carbone dissous, plus les ions carbonate sont consommés : il en reste moins
pour former le carbonate de calcium des coquilles et des squelettes.
\item On applique $\hho=\cst\times10^{-\pH}$ :
\begin{gather*}
\hho_{1985}=10^{-8{,}11}\times\cst=\res{\qty{7,8e-9}{mol.L^{-1}}}\\
\hho_{2024}=10^{-8{,}04}\times\cst=\res{\qty{9,1e-9}{mol.L^{-1}}}
\end{gather*}
Le rapport vaut $10^{0{,}07}\approx1{,}175$ : l'augmentation relative est d'environ \res{18~\%}.
\item Le pH de 2100 serait compris entre $8{,}04-0{,}50=7{,}54$ et $8{,}04-0{,}15=7{,}89$, d'où
\[\res{\qty{1,3e-8}{mol.L^{-1}}\leq\hho\leq\qty{2,9e-8}{mol.L^{-1}}}.\]
Le facteur d'augmentation serait compris entre $10^{0{,}15}\approx1{,}4$ et $10^{0{,}5}\approx3{,}2$.
\item Même en 2100, le pH resterait supérieur à 7 : l'eau de mer resterait basique. « Acidification »
signifie que le pH \emph{diminue}, c'est-à-dire que \hho\ augmente ; il serait plus juste d'écrire :
« les océans deviennent moins basiques ».
\end{questions}
\end{corrige}
\end{blocLivre}
\begin{blocLivre}{cor-1-15}
\begin{corrige}{15}{Tester la relation entre pH et concentration}
\demarche{Choisir une verrerie de précision, comparer les mesures au modèle, puis interpréter chaque écart.}
\begin{questions}
\item La dilution conserve la quantité de soluté : $c_0V_0=c_1V_1$, d'où $V_0=\frac{V_1}{10}=\qty{10,0}{mL}$.
Prélever 10,0 mL de S$_0$ avec une \res{\text{pipette jaugée de 10,0 mL}} munie d'une propipette, les verser
dans une \res{\text{fiole jaugée de 100,0 mL}}, compléter avec de l'eau distillée jusqu'au trait de jauge,
boucher et agiter. La verrerie jaugée garantit des volumes précis.
\item Étalonner le pH-mètre avec les solutions étalons ; rincer la sonde à l'eau distillée et l'essuyer
entre deux mesures ; mesurer de la solution la plus diluée à la plus concentrée. Porter des lunettes et des gants.
\item On attend $a=1$ et $b=0$. On obtient $a\approx0{,}97$ et $b\approx0{,}11$ : les cinq premières mesures
sont \res{\text{compatibles avec la relation}}, à quelques centièmes d'unité près.
\item Pour $c=\qty{1,0e-1}{mol.L^{-1}}$, le pH mesuré (1,10) dépasse la valeur attendue (1,00) : à forte
concentration, les ions interagissent et le modèle de la solution diluée atteint ses limites.
Pour $c=\qty{1,0e-6}{mol.L^{-1}}$, on mesure 5,81 au lieu de 6,00, soit $\hho\approx\qty{1,5e-6}{mol.L^{-1}}$ :
une autre source d'ions oxonium intervient, le dioxyde de carbone de l'air dissous dans l'eau distillée
(\ce{CO2,H2O + H2O <=> HCO3- + H3O+}). Elle est négligeable aux fortes concentrations, pas ici.
\item Les écarts valent 0,10 ; 0,04 ; 0,02 ; 0,01 ; 0,02 et 0,19. La relation est vérifiée à 0,05 près
\res{\text{pour } c \text{ compris entre } \num{1.0e-5} \text{ et } \qty{1,0e-2}{mol.L^{-1}}}.
\end{questions}
\schema[corrige]{01-ex15-graphe}
\end{corrige}
\end{blocLivre}
